% !TEX TS-program = lualatexmk
% !TEX parameter =  --shell-escape
\documentclass[12pt]{book}

%%%%%%%%%%%%%%% STANDARD PACKAGES %%%%%%%%%%%%%%%

%% for personal reasons, the commands are a mix of French and English names. Feel free to choose one or the other if you make use of them. This file is subject to change later on.

\usepackage{geometry}
\usepackage{fontspec}
\usepackage[english,ecclesiasticlatin.usej,activeacute]{babel}
     \babelprovide[hyphenrules=latin]{ecclesiasticlatin}
%          \usepackage{xspace} %% better to use {} after certain macros (or even within definitions) 
     \usepackage{multicol}
\usepackage{paracol}
\footnotelayout{m} %% this allows merged footnote if text is in parallel columns (mostly translations)
\usepackage{fancyhdr}
\usepackage{needspace} %%pour réserver de l'espace en bas de page ; ça évite des séparations entre les titres et la partition ou le texte qui les suivent.
\usepackage[compact]{titlesec}
\usepackage{emptypage}
\usepackage{xcolor}
\usepackage{perpage}%%%to reset footnotes at each page%%%
%\usepackage{xstring}
\usepackage{enumitem} %% nécessaire pour les textes des psaumes.
\usepackage{etoolbox}
\usepackage{lettrine}
%\usepackage{tocloft}%%to fix toc title. titlesec companion packages could also work.
\usepackage{indentfirst}%% to allow for \section with indented 1st paragraphs in good Franco-Latin style (LaTeX default is to suppress this. Package sets  \@afterindentfalse to (always) true.
\usepackage[savepos]{zref}
\usepackage{hyperref}
\usepackage{refcount}


%%%%%%%%%%%%%%% HYPHENATION AND TYPOGRAPHICAL CONVENTIONS %%%%%%%%%%%%%%

%% page settings %%%%
\sloppy %% to avoid overfull hboxes

%\setlength{\textwidth}{5.75in}

\geometry{letterpaper}
\geometry{bindingoffset=0.5in,
inner=1in,
outer=1.35in,
top=1.25in,
bottom=1.25in,
headsep=0.75in
} %%

%% General scale of all graphical elements. Also borrowed from MB
%% Values different from 1 are largely untested.
%% Used in those commands (e.g. everything FontSpec) that use a scale parameter.
\newcommand{\customscale}{1}

\AtBeginDocument{\setlength{\parindent}{1em}} %% should set 1em to EB Garamond not Computer Modern.

%%%%%%%%%%%%%%% GREGORIO CONFIG %%%%%%%%%%%%%%%

\usepackage[autocompile]{gregoriotex}

%% sets * and † to font called via fontspec
\def\GreStar{*}
\def\GreDagger{†}

%%prints asterisk instead of star normally inserted by <v>\greheightstar</v>
 \gredefsymbol{greheightstar}{EB Garamond}{asterisk}

%% should prevent flat turning custos into a clef, as Solesmes  uses an altered custos in only one chant per Élie Roux.
\gresetcustosalteration{invisible}

%produces a score with a smaller initial (default is 40 pt) and annotations like in the Solesmes books; the 1st is optional and can be omitted as needed.
 \NewDocumentCommand{\gscore}{O{}mm}{%
 \greannotation{#1}%
\greannotation{#2}%
\grechangestyle{initial}{\fontsize{28}{28}\selectfont}%
 \gregorioscore{partitions/#3}}
 
 
  \NewDocumentCommand{\lectionscore}{O{}m}{%
   \grecommentary{#1}%
  \grechangestyle{initial}{\fontsize{36}{36}\selectfont}%
  \gregorioscore{partitions/#2}}
 
 %% default should be 40, if not using, annotations need to be turned off in gabc. For the first antiphon of the office, Gospel canticle etc. (initial is both to refer to first and the large illustrated initial, used sometimes with an annotation, cf. Christus factus est of Tenebrae in the ritus monasticus.
  \NewDocumentCommand{\initialscore}{O{}m}{%
      \grechangedim{initialraise}{0.5cm}{scalable}%
  \greannotation{#1}%
   \grechangestyle{initial}{\fontspec{EB Garamond Initials}\fontsize{45}{45}\selectfont}%
 \gregorioscore{partitions/#2}
   \grechangedim{initialraise}{0cm}{scalable}}

%% score with no initial, e.g psalms, common tones, etc.
\newcommand{\smallscore}[2][y]{
  \gresetinitiallines{0}
  \gregorioscore{partitions/#2} %% à remplacer avec NewDocumentCommand ; les arguments ne marchent pas correctement … que fait le [y] ?
  \gresetinitiallines{1}
}

\NewDocumentCommand{\zerobaroffsettextleft}{}%
{%
\grechangedim{maxbaroffsettextleft}%
{0cm}%
{scalable}%
}

\NewDocumentCommand{\zerobaroffsettextright}{}%
{%
\grechangedim{maxbaroffsettextright}%
{0cm}%
{scalable}%
}

%%%the default for left bar offset is 0.3cm. Right is 0.15cm. This needs to be applied after a score where the bar offset is modified.

\NewDocumentCommand{\resetbaroffsettextleft}{}%
{%
\grechangedim{maxbaroffsettextleft}%
{0.3cm}%
{scalable}%
}

\NewDocumentCommand{\resetbaroffsettextright}{}%
{%
\grechangedim{maxbaroffsettextright}%
{0.15cm}%
{scalable}%
}

 \NewDocumentCommand{\MagAnnotation}{}{%
 \grechangedim{annotationseparation}{-0.01cm}{fixed}%
 } %% pour les antiennes où le chiffre convient mieux aux miniscules (parfois 1, 2, 4, 7…) ; il n'est pas nécessaire pour 6 ou 8.
%\gresetheadercapture{annotation}{greannotation}{string}

%% A-bove L-ines T-ext shall be italicized and in a smaller font
\grechangestyle{abovelinestext}{\itshape\fontsize{8}{10}\selectfont}
%% for psalm tones etc. Roman/upright text is what the Liber Usualis uses. 
\NewDocumentCommand{\altnormal}{}{\grechangestyle{abovelinestext}{\normalfont\fontsize{8}{10}\selectfont}}
%%  resets ALT font
\NewDocumentCommand{\altitshape}{}{\grechangestyle{abovelinestext}{\itshape\fontsize{8}{10}\selectfont}} 

\NewDocumentCommand{\altlarge}{}{\grechangestyle{abovelinestext}{\large}} %% for I reading of Lamentations

\NewDocumentCommand{\altupright}{}{\grechangestyle{abovelinestext}{\normalfont}} %%


%% fine-tuning of space beween the staff and the text above lines (used for Magnificats; needs to be rethought (again, should 2 and 3 be raised or not?) and worked into \smallscore
\newcommand{\altraise}{0.1cm} %% default is -0.1cm %% needs to be fine-tuned for \rubrique command with EB Garamond font. -0.2cm is MB value
\grechangedim{abovelinestextraise}{\altraise}{scalable}

\newcommand{\altheight}{0.5cm} %% default is 0.3cm and value must be bigger than text height; this should fix the problem. but needs further investigation.
\grechangedim{abovelinestextheight}{\altheight}{scalable} %% this is a very finicky command.

\newcommand{\comraise}{0.4cm} %% default is 0,2m %%
\grechangedim{commentaryraise}{\comraise}{scalable}

%% allows printing of glyphs from Gregorio score font (available glyphs listed in documentation) as text.
\makeatletter
\def\gretextglyph#1{{\gre@font@music\csname GreCP#1\endcsname}}
\makeatother

%%% Font %%%
\setmainfont{EB Garamond}[UprightFont=EB Garamond Regular,
ItalicFont= EB Garamond Italic,
BoldFont= EB Garamond Bold,
Ligatures=Rare,
Numbers=Proportional,
Numbers=OldStyle] %% all \kern numbers are based on this and can be removed if this font is abandoned.
\newfontfamily\symbolfont{liturgy}
  %% text must be written {\symbolfont{+}} (you can use V and R as well) but may conflict as + is † in gabc. Cross should be \small or even \footnotesize


%    \grechangestyle{initial}{\fontspec{EB Garamond Initials}\fontsize{#1}{#1}\selectfont}

%%% default initial size is 32 points
%\newcommand{\defaultinitialsize}{28}
%\initialsize{\defaultinitialsize}

%%%% Font %%%
%\setmainfont{EB Garamond}[UprightFont=EB Garamond Regular,
%ItalicFont= EB Garamond Italic,
%BoldFont= EB Garamond Bold,
%Ligatures=Rare,
%StylisticSet=6,
%Numbers=Proportional,
%Numbers=OldStyle] %% all \kern numbers are based on this and can be removed if this font is abandoned. %%StylisticSet=6 is, in Pardo EBGaramond, the long-tailed Q.
%\newfontfamily\symbolfont{liturgy} 
  %% text must be written {\symbolfont{+}} (you can use V and R as well) but may conflict as + is † in gabc. Cross should be \small or even \footnotesize
  
  %%%Roman numerals for the year %% https://tex.stackexchange.com/questions/185548/the-year-in-roman-and-the-month-in-text

\makeatletter
\newcommand{\YEAR}{\@Roman{\the\year}}
\makeatother


%%% TYPOGRAPHICAL AND SYMBOL COMMANDS  %%%%

%% \P is pilcrow and is defined in LaTeX as is.
%% Add {\textcolor{gregoriocolor} as first part of command's definition if changing to red.


%%for proper spacing of all-caps letters to prevent clashes, e.g. serif strokes touching.

\NewDocumentCommand\capspace{m}{{\addfontfeature{LetterSpace=5.0}{#1}}}
\NewDocumentCommand\scspace{m}{\textsc{{{\addfontfeature{LetterSpace=5.0}{#1}}}}}

\NewDocumentCommand\feasttitle{m}{{\centering{\capspace{\huge{#1}}}\par}}


%% macro to print Alleluia for versicles.
\NewDocumentCommand{\tpalleluia}{}{(\textit{T.P.} Allelúia.)}

\newcommand{\specialcharhsep}{3mm} % space after invoking R/ or V/ or A/ outside rubrics
%\newcommand{\vv}{\textcolor{gregoriocolor}{\fontspec[Scale=\customscale]{Charis SIL}℣.\nolinebreak[4]\hspace{\specialcharhsep}\nolinebreak[4]}} %% format from MB

\newcommand{\vv}{{\normalfont ℣.\nolinebreak[4]\hspace{\specialcharhsep}\nolinebreak[4]}}
\newcommand{\rr}{{\normalfont ℟.\nolinebreak[4]\hspace{\specialcharhsep}\nolinebreak[4]}}

%% typesets a cross pattée.
\NewDocumentCommand{\cc}{}{%
    % Grouping to keep font changes local
    {%
        % Ensure we're not in italics (since liturgy.ttf doesn't have italic)
        \normalfont%
        % Select the font liturgy.ttf
        \symbolfont%
        % Set the size
        \footnotesize%
        % In liturgy.ttf, the plus (+) is a cross pattée
        +%
    }%    
} %%  can replace <+> with <U+2720> (see above) if a suitable font is found (or EB Garamond font is fixed); command will be useful in lieu of typing the Unicode.
%% use <v> tag to insert into a gabc score (usually for the faithful or for blessings, e.g. the font.

%% Same special characters, for in-score use (<sp>V/ R/ A/ +</sp>)

\gresetspecial{V/}{{\fontspec[Scale=\customscale]{EB Garamond}℣.~}}
\gresetspecial{R/}{{\fontspec[Scale=\customscale]{EB Garamond}℟.~}}
\gresetspecial{+}{{\fontspec[Scale=\customscale]{EB Garamond}†~}}
\gresetspecial{*}{\gresixstar}
\gresetspecial{cross}{\cc}

%% Same special characters, for use in rubrics (no space)

\newcommand{\vvrub}{{\normalfont ℣.~}}
\newcommand{\rrrub}{{\normalfont ℟.~}}

\NewDocumentCommand{\antiphona}{mm}
{\begin{otherlanguage*}{#1}%
\noindent%
 Ant. #2.%
 \end{otherlanguage*}%
 }
 


%%rubrics: black italics, smaller than body of psalms etc

\NewDocumentCommand{\rubrique}{m}{%
    {%
        \fontsize{10}{12}%
        \selectfont%
        \textit{%
        %
        #1%
    }%
    }}
    
%% macro to print normal text inside of rubric (name of a chant or prayer, etc.)

\NewDocumentCommand{\normaltext}{m}{%
    {%
        \normalfont%
        #1%
    }%
    }
    
%% in case something should be bolded inside of a rubrique
\NewDocumentCommand{\rubriquegras}{m}{%
    {%
        \normalfont%
\textbf{#1%
}%
}}

%% to print in red instead of italicizing
\NewDocumentCommand{\rouge}{m}{%
\textcolor{gregoriocolor}%
{\fontsize{10}{12}%
\selectfont{%
#1%
}%
}}

%% to print in black within rouge group.

%\NewDocumentCommand{\textenoir}{m}{%

\NewDocumentCommand{\textenoir}{m}{%
\textcolor{black}%
{%
\normalfont%
#1%
}%
}

%% rouge but in italic (this is technically not correct).
\NewDocumentCommand{\rougeit}{m}{%
\textcolor{gregoriocolor}%
{\fontsize{10}{12}%
\selectfont%
\textit{#1%
}%
}}

%%for Scriptural references of the chapter.
\NewDocumentCommand{\scripture}{m}{%
{\raggedleft{\textit{#1}%
}%
\par}%
}

%\newcommand{\rubriquegras}[1]{{\fontsize{9}{11}\selectfont\textbf{#1}}}

%%macro to space punctuation with ecclesiasticlatin language from babel.
\NewDocumentCommand{\espaceponctuation}{}{\hspace{0.10em}} %%this value seems more balanced with this font than the definition provided in the ecclesiasticlatin documentation.

\NewDocumentCommand{\myrule}{}{%
    \par%
    {%
        \centering%
        \rule{0.3\textwidth}{0.4pt}%
        \par%
    }% typesets a horizontal rule like on the page with the prayers before and after the office.
}  %%If you're using color at all in your document, you might want to either force \myrule to use black or make it cusotmizable. (from u/Independent-Comb-257)

\NewDocumentCommand{\bigrule}{}{%
    \par%
    {%
        \centering%
        \rule{0.6\textwidth}{0.4pt}%
        \par%
    }% typesets a horizontal rule like on the page with the prayers before and after the office.
}  %%If you're using color at all in your document, you might want to either force \myrule to use black or make it cusotmizable. (from u/Independent-Comb-257)


%%% typesets a cross pattée.

\NewDocumentCommand{\mycross}{}{%
{\symbolfont\small{{+}}}%
{}
} %% can replace <+> with <U+2720> if a suitable font is found (or EB Garamond font is fixed); command will be useful in lieu of typing the Unicode.

%% macro to format psalm text

%\NewDocumentCommand{\pstexte}{ m }{%
%    \smallskip%
%    \noindent%
%    \begin{itemize}[%
%    		label=\null, %
%			leftmargin=0pt, %
%			itemindent=0mm, %
%			labelsep=0pt, %
%			labelwidth=0pt, %
%			rightmargin=0pt, %
%			parsep=0pt, %
%			topsep=0pt, %
%			itemsep=0pt]%
%    \input{psaumes/#1}
%    \end{itemize}}
          
      \setlength{\columnsep}{3pc}
%\setlength{\columnsep}{10pt}
\setlength\columnseprule{0.4pt}


%% original version kept so that I can fix errors found later and keep using common texts, going forward
    \NewDocumentCommand{\textebilingue}{ mm }{%
    \smallskip%
       \begin{paracol}{2}
           \input{psaumes/#1}
           \switchcolumn
    \input{psaumes/#2}
    \end{paracol}}

%2nd version;  no need for input (too many chapter/collect files) + is the equivalent of long, allows \par
%    \NewDocumentCommand{\texteparacol}{ +m+m }{%
%    \smallskip%
%       \begin{paracol}{2}
%           #1
%           \switchcolumn
%    #2
%    \end{paracol}}

%%%3rd version: need to switch languages!
    \NewDocumentCommand{\texteparacol}{ +mm+m }{%
    \smallskip%
       \begin{paracol}{2}
           #1
           \switchcolumn
           \begin{otherlanguage}{#2}
    #3
    \end{otherlanguage}
    \end{paracol}}
    
    
    %%psalm vernacular file has language command
    \NewDocumentCommand{\pstexte}{mm}{%
    \smallskip%
    \noindent%%
   \begin{paracol}{2}
    \begin{enumerate}[%
    		label=\arabic*., %
			leftmargin=3pt, %
			itemindent=3mm, %
			labelsep=3pt, %
			labelwidth=0pt, %
			rightmargin=0pt, %
			parsep=0pt, %
			topsep=0pt, %
			itemsep=0pt,%
			start=2]%
    \input{psaumes/#1.tex}
    \end{enumerate}
    \switchcolumn
       \begin{enumerate}[%
    		label=\arabic*.,%
			leftmargin=3pt, %
			itemindent=3mm, %
			labelsep=3pt, %
			labelwidth=0pt, %
			rightmargin=0pt, %
			parsep=0pt, %
			topsep=0pt, %
			itemsep=0pt]%
    \input{psaumes/#2.tex}
      \end{enumerate}
    \end{paracol}%
    } %% modified from original in order to allow 2 column psalms
    
    %% Command de Matthias Bry modifié, prints psalm incipit score with the text
    \NewDocumentCommand{\psalmus}{mmmm}{%
	\needspace{4\baselineskip}%
	\smalltitle{Psalmus #1.}%
	\smallscore[n]{#2}%
	\pstexte{#3}{#4}%
}

    %% sd=semiduplex. especially for Dominica ad Vesperas in the Psalterium (the only example which comes to mind, in fact!)
 \NewDocumentCommand{\sdpsalmusbilingue}{mmm}{%
	\needspace{4\baselineskip}%
		\smalltitle{Psalmus #1.}%
	\pstexte{#2}{#3}%
	}

    %% Command de Matthias Bry modifié, prints canticle incipit score with the text and scriptural commentary
    \NewDocumentCommand{\canticum}{mmmm}{
	\needspace{4\baselineskip}
	\smalltitle{Canticum #1.}
	 \grecommentary{#2}
	\smallscore[n]{#3}
	\pstexte{#4}
}
	
    %% macro to print any additional text (Capitiulum, oratio, rubrics)
 \NewDocumentCommand{\textes}{ m }{%
    \input{textes/#1}}
    
%     \NewDocumentCommand{\lection}{m}{%
%   \begin{multicols}{2}#1\end{multicols}} %% this doesn't work as expected, needs to be fixed
    
    
    
    %% SC work for page headers and for secondary headings but not so much for things like "Dominica…" and certainly not the feast name%%
    
    %%%% Headers%%%%
    \pagestyle{fancy}
\fancyhead{}
\fancyfoot{}
\renewcommand{\headrulewidth}{0pt}


 \setlength{\headheight}{13.59999pt} %% recommended by LaTeX

%\fancyhead[RO]{\small\thepage}
%\fancyhead[LE]{\small\thepage}

\fancyhead[RO]{\small\rightmark\hspace{0.5cm}\thepage}
\fancyhead[LE]{\small\thepage\hspace{0.5cm}\leftmark}

\newcommand{\setheaders}[2]{
	\renewcommand{\rightmark}{{\sc#2}}
	\renewcommand{\leftmark}{{\sc#1}}
}
\setheaders{}{}


%% TITLE Commands %%%% %% TITLE Commands %%%% currently dead but Matthias says this isn't a good idea.
 \titleformat{\chapter}[display]{\Huge\filcenter\sc\addfontfeature{LetterSpace=5.0}}{}{0pt}{} %%
\titleformat{\section}[display]{\LARGE\filcenter\sc\addfontfeature{LetterSpace=5.0}}{}{}{} %%
\titleformat{\subsection}[display]{\Large\filcenter\sc\addfontfeature{LetterSpace=5.0}}{}{}{} %%originally \large
\setcounter{secnumdepth}{0}
\titlespacing*{\chapter}{0pt}{-25pt}{20pt} %%this should reduce spacing before chapter title while putting it flush with the top of the text area
%% see https://tex.stackexchange.com/questions/63390/how-to-decrease-spacing-before-chapter-title
%%\titlespacing*{\chapter}{0pt}{-50pt}{40pt}  %%this puts chapter title in HEADER which we do not want
\titlespacing{\section}{0pt}{*0}{*0}

%%need section titleformat more appropriate for weekdays of psalter and minor feasts 
% See https://tex.stackexchange.com/questions/623797/multiple-section-styles-in-same-document.

\addto\captionsecclesiasticlatin{\renewcommand\contentsname{\centerline{INDEX GENERALIS.}}}

% \NewDocumentCommand{\mysection}{m}{\section{#1}\setheaders{{\scshape\addfontfeature{LetterSpace=5.0}#1}}{{\scshape\addfontfeature{LetterSpace=5.0} #1}}}
 
\NewDocumentCommand{\header}{m}{\setheaders{{\scshape\addfontfeature{LetterSpace=5.0}#1}}{{\scshape\addfontfeature{LetterSpace=5.0} #1}}}

%% all of these are a mess and should probably be ignored, but feel free to use them by referring to Matthias B's originals in the Nocturnale Romanum repository.

\newcommand{\officiumtitulum}[1]{
%  \newpage
%  \thispagestyle{empty}
  \setheaders{{\scshape\addfontfeature{LetterSpace=3.0}#1}}{{\scshape\addfontfeature{LetterSpace=3.0} #1}}
  \begin{center}
  {\scshape\large #1}\par
  \end{center}}
  
%   \setheaders{{{\scshape\addfontfeature{LetterSpace=3.0}}}{{{\scshape\addfontfeature{LetterSpace=3.0} {#1}}}

%%formatting date of feasts where this listed at top of page
\NewDocumentCommand{\litdate}{m}{%
\smalltitle{Die #1.}%
}

\NewDocumentCommand{\litdateen}{m}{%
\smalltitle{#1.}%
}

%%formatting for Sunday feasts : inter, post… name is for consistency
\NewDocumentCommand{\dominicadate}{m}{%
\smalltitle{Dominica #1.}%
}

\NewDocumentCommand{\sundaydate}{m}{%
\smalltitle{Sunday #1.}%
}

%%will use litdaten for English
%%formatting for feasts that fall on a weekday fixed to some other thing, like Corpus Christi, Sacred Heart, BVM in Passiontide, and the new feasts found in the appendix (Feria Roman numeral in Latin: inter, post… name is for consistency
\NewDocumentCommand{\feriadate}{m}{%
\smalltitle{Feria #1.}%
}

 %%looks like xparse formatting changed so o just has No Default Value whereas O{} is that OR you can insert something in braces!
%%answer thanks to David Carlisle
\NewDocumentCommand{\duplexclassis}{mo}{%
\needspace{5\baselineskip}%
\vspace{\baselineskip}%
 {\centering\textit{\large Duplex #1 classis\IfValueT{#2}{ #2}.}\par}%
}

\NewDocumentCommand{\doubleclass}{mo}{%
\needspace{5\baselineskip}%
\vspace{\baselineskip}%
 {\centering\textit{\large Double of the #1 class\IfValueT{#2}{ #2}.}\par}%
}


\NewDocumentCommand{\duplexmajus}{}{%
\needspace{5\baselineskip}%
\vspace{\baselineskip}%
 {\centering{\textit{\large Duplex majus.}}\par}%
}

\NewDocumentCommand{\doublemajor}{}{%
\needspace{5\baselineskip}%
\vspace{\baselineskip}%
 {\centering{\textit{\large Double major.}}\par}%
}

%%is Duplex even specified.as it's the default?
\NewDocumentCommand{\duplex}{}{%
\needspace{5\baselineskip}%
\vspace{\baselineskip}%
 {\centering{\textit{\large Duplex.}}\par}%
}

\NewDocumentCommand{\double}{}{%
\needspace{5\baselineskip}%
\vspace{\baselineskip}%
 {\centering{\textit{\large Double.}}\par}%
}

\NewDocumentCommand{\semiduplex}{}{%
\needspace{5\baselineskip}%
\vspace{\baselineskip}%
 {\centering{\textit{\large semiuplex.}}\par}%
}

\NewDocumentCommand{\semidouble}{}{%
\needspace{5\baselineskip}%
\vspace{\baselineskip}%
 {\centering{\textit{\large Semidouble.}}\par}%
}

\NewDocumentCommand{\simplex}{}{%
\needspace{5\baselineskip}%
\vspace{\baselineskip}%
 {\centering{\textit{\large Simplex.}}\par}%
}

\NewDocumentCommand{\simple}{}{%
\needspace{5\baselineskip}%
\vspace{\baselineskip}%
 {\centering{\textit{\large Simple.}}\par}%
}

\NewDocumentCommand{\rank}{m}{
\needspace{5\baselineskip}  %% very small if there are neumes above the staff, including flats in mode 2, e.g. O Doctor optime
\vspace{\baselineskip}
 {\centering{\textit{\large #1}}\par}
}



\NewDocumentCommand{\bigtitle}{m}{
\needspace{5\baselineskip}  %% very small if there are neumes above the staff, including flats in mode 2, e.g. O Doctor optime
\vspace{\baselineskip}
 {\centering{\large#1}\par}
}

\NewDocumentCommand{\smalltitle}{m}{
\needspace{5\baselineskip}  %% very small if there are neumes above the staff, including flats in mode 2, e.g. O Doctor optime
\vspace{\baselineskip}
 {\centering #1\par}
}

\NewDocumentCommand{\chapterlateng}{}{%
\smalltitle{Chapter.}%
}


%\newcommand{\subtitulum}[1]{
% \subsection{
%  \begin{center}
%  {\large#1}\par
%  \end{center}}}
  
  %% something has broken here
  %% see https://tex.stackexchange.com/questions/545841/paragraph-ended-before-ttlformatsi-was-complete

%% originally used apparently for typesetting portions of the Common of the BVM
\newcommand{\espacetitre}{\vspace{-0.5cm}}
\newcommand{\espaceps}{\vspace{-3mm}} %% ps =psaume
\newcommand{\espacecap}{\vspace{-3mm}} %% cap is capitulum
\newcommand{\espace}{\vspace{2mm}}

\NewDocumentCommand{\secreto}{}{Pater noster. \rubrique{secreto.}}

\NewDocumentCommand{\pateravecredo}{}{\large{Pater noster, Ave María, \textit{et} Credo.}}

\NewDocumentCommand{\pateravedeus}{}{%
Pater noster et Ave María. ℣. Deus in adjutórium.%
}

\NewDocumentCommand{\prayers}{}{%
Our Father and Hail Mary, \rubrique{silently.}%
}

\NewDocumentCommand{\orationes}{}{%
Pater noster et Ave María \rubrique{secreto.}%
}

%%for litanies
\makeatletter  
\newcounter{score}
\newcounter{tabstop}[score]
\newcommand{\grealign}{%
	\@bsphack%
	\ifgre@boxing\else%
		\kern\gre@dimen@begindifference%
		\stepcounter{tabstop}%
		\expandafter\zsavepos{stop-\thescore-\thetabstop}%
		\kern-\gre@dimen@begindifference%
	\fi%
	\@esphack%
}

\newcommand{\setstops}{%
  \gdef\nstabbing@stops{%
    \hspace*{-\oddsidemargin}\hspace{-1in}%
    \hspace*{\zposx{stop-\thescore-1} sp}\=%
  }%
  \count@=\@ne
  \loop\ifnum\count@<\value{tabstop}%
    \begingroup\edef\x{\endgroup
      \noexpand\g@addto@macro\noexpand\nstabbing@stops{%
        \noexpand\hspace{-\noexpand\zposx{stop-\thescore-\the\count@} sp}%
        \noexpand\hspace{\noexpand\zposx{stop-\thescore-\the\numexpr\count@+1} sp}\noexpand\=%
      }%
    }\x
    \advance\count@\@ne
  \repeat
  \nstabbing@stops\kill
}
\makeatother

\newenvironment{nstabbing}
  {\setlength{\topsep}{0pt}%
   \setlength{\partopsep}{0pt}%
   \tabbing%
   \setstops}
  {\endtabbing\stepcounter{score}}

\makeatletter
\newcounter{blankpages} %% should remove number on blank page at end of preface. Cf https://tex.stackexchange.com/questions/250761/how-to-not-count-blank-pages-in-frontmatter-of-two-sided-document
\def\cleardoublepage{%
    \clearpage
    \if@twoside
    \ifodd\c@page
    \else
    \hbox{}
    \thispagestyle{empty}%
    \newpage\stepcounter{blankpages}%
    \if@twocolumn\hbox{}\newpage\fi
    \fi
    \fi
}
\newcommand{\@romannoblank}[1]{%
    \@roman{\numexpr#1-\value{blankpages}\relax}%
}
\makeatother

\providecommand\phantomsection{} %% should make hyperref work properly %%only works with section etc., not with label etc. (that requires \phantomsection
%

\NewDocumentEnvironment{enpars}{}
  {\begin{otherlanguage*}{english}}
  {\par\end{otherlanguage*}}
  
  \NewDocumentEnvironment{frpars}{}
  {\begin{otherlanguage*}{french}}
  {\par\end{otherlanguage*}}
%\renewcommand{\contentsname}{\hfill\bfseries\LARGE Index Generalis.\hfill\null}  %% modifies TOC title
%\renewcommand{\cfttoctitlefont}{\hfil\LARGE}

%%%SUBFILES%%%
\usepackage{xr} %%to allow cross-references between documents%%%
  \usepackage{subfiles}
  
  %% When we start a new subfile (new chapter or section) , 
%% we start on a new page (with blank filling on the previous page) and create a corresponding label.

\newcommand{\customsubfile}[1]{\newpage\label{#1}\thispagestyle{empty}\subfile{#1}}


\begin{document}
 
 \twosided[pc] %% to allow for the left-right switching of columns needed for bilingual liturgical typesetting (like in missals); it's unclear why this goes BEFORE everything else and not when activating the environment. %% It's OK for paper booklets to not switch, but it's probably best practice to switch.
 
 \thispagestyle{empty}
 
 {\centering{\huge\capspace{{SUNDAYS OF PASCHAL TIME}}}\par}\header{sundays of paschal time at vespers.}
   \vspace{0.5\baselineskip}
 
  {\centering{\huge\capspace{{AT VESPERS.}}}\par}
  
\smalltitle{Before the Office.}

\rubrique{These prayers are said kneeling. End when the priest stands.}

\textebilingue{orationes_ante_do}{orationes_ante_do_eng}

\smalltitle{\prayers}

\rubrique{The celebrant intones the following verse. All sing the response and doxology.}

\begin{enpars}
\rubrique{Make the sign of the cross at the beginning and at the Magnificat, and bow the head when the Trinity is mentioned, when mention is made of the ``Holy Name'' in some way, and at ``Orémus.''}\end{enpars}

\gscore[]{℣.}{Deus_In_Adjutorium_Festivus}

\noindent
\begin{enpars}
O God, \cc{} come to my assistance. \rr O Lord, make haste to help me. Glory be to the Father, and to the Son,~* and to the Holy Ghost. As it was in the beginning, is now,~* and ever shall be, world without end. Amen. Alleluia.\end{enpars}

\smalltitle{Antiphon. 7. c2.}
\gscore[]{}{alleluia_sun_tp_vespers}

\rubrique{This antiphon is sung over all of the psalms. On the I Sunday after Easter, the office is of double major rite, so the antiphon \normaltext{Allelúia, \pageref{duplex_antiphon}} is sung before and after.}

\rubrique{\begin{enpars}
Sit at the word ``Sede.'' The cantor sings odd verses; all sing even verses, and this for the hymn and Magnificat as well.
\end{enpars}}

\psalmus{109}{109_7c2}{109_7}{109_en}

\rubrique{\begin{enpars}
The cantors then begin the psalms at the conclusion of the doxology.
\end{enpars}}

\psalmus{110}{110_7c2}{110_7}{110_en}

\psalmus{111}{111_7c2}{111_7}{111_en}

\psalmus{112}{112_7c2}{112_7}{112_en}

\psalmus{113}{113_7c2}{113_7}{113_en}

\label{duplex_antiphon}

\smallscore{an_alleluia_sun_tp_vespers_solesmes_1961}
\antiphona{english}{Alleluia, alleluia, alleluia}

\rubrique{\begin{enpars}
Having sung the last antiphon in full, stand when the celebrant does.
\end{enpars}}

%\newpage

\chapterlateng

\smalltitle{I Sunday after Easter.}

\scripture{1 Joann. 5, 4.}
\texteparacol{\lettrine{C}{a}ríssimi: Omne quod natum est ex Deo, vincit mundum:~† et hæc est victória, quæ vincit mundum,~* fides nostra.}{english}{\noindent My beloved: For whatsoever is born of God, overcometh the world: and this is the victory which overcometh the world, our faith.}

\smalltitle{II Sunday after Easter.}

\scripture{1 Petri 2, 21 – 22.}

\texteparacol{\lettrine{C}{a}ríssimi: Christus passus est pro nobis,~† vobis relínquens exémplum ut sequámini vestígia ejus.~* Qui peccátum non fecit, nec invéntus est dolus in ore ejus.}{english}{\noindent Dearly beloved: Christ also suffered for us, leaving you an example that you should follow his steps. Who did no sin, neither was guile found in his mouth.}

\smalltitle{III Sunday after Easter.}

\scripture{1 Petri 2, 11.}
\texteparacol{\lettrine{C}{a}ríssimi: Obsecro vos tamquam ádvenas et peregrínos,~† abstinére vos a carnálibus desidériis,~* quæ mílitant advérsus ánimam.}{english}{\noindent Dearly beloved, I beseech you as strangers and pilgrims, to refrain yourselves from carnal desires which war against the soul.}

\smalltitle{IV Sunday after Easter.}

\scripture{Jac. 1, 17.}
\texteparacol{\lettrine{C}{a}ríssimi: Omne datum óptimum, et omne donum perféctum desúrsum est, descéndens a Patre lúminum,~† apud quem non est transmutátio,~* nec vicissitúdinis obumbrátio.}{english}{\noindent Dearly beloved: Every best gift, and every perfect gift, is from above, coming down from the Father of lights, with whom there is no change, nor shadow of alteration.}

\smalltitle{V Sunday after Easter.}

\scripture{Jac. 1, 22 – 24.}

\texteparacol{\lettrine{C}{a}ríssimi: Estóte factóres verbi, et non auditóres tantum: falléntes vosmetípsos.~† Quia si quis audítor est verbi, et non factor: hic comparábitur viro consideránti vultum nativitátis suæ in spéculo:~* considerávit enim se, et ábiit, et statim oblítus est qualis fúerit.}{english}{\noindent Dearly beloved: But be ye doers of the word, and not hearers only, deceiving your own selves. For if a man be a hearer of the word, and not a doer, he shall be compared to a man beholding his own countenance in a glass. For he beheld himself, and went his way, and presently forgot what manner of man he was.}

%% The usual tone is  omitted as we do not use these tones, but they are included in the tex file and the gabc files are available, for the sake of completeness.

\hymnlateng

%\gscore[]{8.}{hy_ad_coenam_agni_providi_textus_antiquus}

%\smalltitle{Another tone ad lib.}

%\grechangenextscorelinedim{3}{spacebeneathtext}{0.25cm}{scalable}
%
%\grechangenextscorelinedim{4,7,10}{spacebeneathtext}{0.03cm}{scalable}

\gscore[]{4.}{hy_ad_coenam_agni_providi_alter_tonus_textus_antiquus}

\textes{ad_coenam_agni_en}

\textes{cantor_rubric}

\smallscore{versiculus_tp}

\begin{enpars} \noindent \vv  Stay with us, o Lord, alleluia.

\noindent \rr  Because it is towards evening, alleluia.
 \end{enpars}
 
 \rubrique{\begin{enpars} 
The celebrant intones the \normaltext{Magnificat} antiphon; this is always proper to the day. The cantor continues the \normaltext{Ma\-gnificat.}\end{enpars}}

\newpage

\bigtitle{Canticle of the Blessed Virgin Mary.}

\smalltitle{I Sunday after Easter.}\label{1_sunday}

\gscore[]{Ant. 8. c}{an_post_dies_octo_solesmes_1961}
\antiphona{english}{After eight days, the doors being shut, the Lord being entered in, said unto them: Peace be unto you, alleluia}

 \smallscore{Magnificat_8c}
 
\magnificattexte{Magnificat_8}

\smalltitle{II Sunday after Easter.}

\gscore[Ad Magnif.]{Ant. 3. a}{an_ego_sum_pastor_bonus_solesmes_1961}
\antiphona{english}{I am the Good Shepherd,  who feed My sheep: and I lay down my life for my sheep. Alleluia}

\smallscore{Magnificat_3a}

\magnificattexte{Magnificat_3a_b}

\smalltitle{III Sunday after Easter.}

\gscore[Ad Magnif.]{Ant. 8. G}{an_amen_amen_quia_plorabitis_solesmes}
\antiphona{english}{Amen, amen I say to you that you shall lament and weep, but the world shall rejoice; and you shall be made sorrowful, but your sorrow shall be turned into joy, alleluia}

 \smallscore{Magnificat_8G}
 
\magnificattexte{Magnificat_8}

\smalltitle{IV Sunday after Easter.}

\gscore[Ad Magnif.]{Ant. 2. D}{an_vado_ad_eum_sed_quia_solesmes_1961}
\antiphona{english}{I go My way to him That sent Me, but because I have said this to you, sadness has filled your hearts, alleluia}

\smallscore{Magnificat_2}

\magnificattexte{Magnificat_2}
\smalltitle{V Sunday after Easter.}

\gscore[Ad Magnif.]{Ant. 8. G*}{an_petite_et_accipietis_solesmes_1961}
\antiphona{english}{Ask, and ye shall receive, that your joy may be full; for the Father himself loveth you, because ye have loved Me, and have believed in Me, alleluia}

 \smallscore{Magnificat_8Gstar}
 
\magnificattexte{Magnificat_8}

\scripture{Luc. 1, 46–55.}

\textes{magnificat_la_en}

 \rubrique{\begin{enpars}
Sit as the antiphon is sung in full. Then the usual dialogue follows. On Sundays, all stand from now until the end of the office, unless one sits for the antiphons sung for the commemorations, according to custom.\end{enpars}}

\textes{dv_la_en}

\myrule

\rubrique{When the officiant is not a deacon:}

\smallscore{domine_exaudi_solemn}

\begin{enpars}
 \vv O Lord hear my prayer.

 \rr And let my cry come unto thee.
\end{enpars}

\myrule
 \rubrique{\begin{enpars}
The celebrant sings the collect of the day, and any commemorations follow. When required, the Suffrage of All Saints is sung in the last place, as below.\end{enpars}}

\bigtitle{Collect.}
\smallskip

\smalltitle{I Sunday after Easter.}

\texteparacol{\lettrine{P}{r}æsta, quǽsumus omnípotens Deus:~† ut qui paschália festa perégimus,~* hæc, te largiénte, móribus et vita teneámus. Per Dóminum.}{english}
{\noindent Grant, we beseech thee, almighty God, that we who have celebrated the paschal feast, may, by thy bounty, retain its fruits in our daily habits and behaviour.
Through our Lord.}
    
\smalltitle{II Sunday after Easter.}

\texteparacol{\lettrine{D}{e}us, qui in Fílii tui humilitáte jacéntem mundum erexísti:~† fidélibus tuis perpétuam concéde lætítiam;~* ut, quos perpétuæ mortis eripuísti cásibus, gáudiis fácias pérfrui sempitérnis. Per eúmdem Dóminum.}{english}
{\noindent O God, whose Son hath humbled himself, and who hast through him raised up the whole world, grant to thy faithful people everlasting joy; and as thou hast delivered them from the bitter pains of eternal death, make them to be glad for ever in thy presence.
Through the same Lord.}
    
    \smalltitle{III Sunday after Easter.}

\texteparacol{\lettrine{D}{e}us, qui errántibus, ut in viam possint redíre justítiæ, veritátis tuæ lumen osténdis:~† da cun\-ctis qui christiána professióne censéntur, et illa respúere quæ huic inimíca sunt nómini;~* et ea quæ sunt apta sectári.
Per Dóminum nostrum.}{english}
{\noindent O God, who showest to them that be in error the light of thy truth, to the intent that they may return into the way of righteousness grant unto all them that are admitted into the fellowship of Christ's religion, that they may eschew those things that are contrary to their profession, and follow all such things as are agreeable to the same.
Through our Lord.}
    
    \smalltitle{IV Sunday after Easter.}

\texteparacol{\lettrine{D}{e}us, qui fidélium mentes uníus éfficis voluntátis:~†  da pópulis tuis id amáre quod prǽcipis, id desideráre quod promíttis;~* ut inter mundánas varietátes, ibi nostra fixa sint corda, ubi vera sunt gáudia. Per Dóminum nostrum.}{english}
{\noindent O God, of whom it cometh that the minds of thy faithful people be all of one will, grant unto the same thy people that they may love the thing which thou commandest, and desire that which thou dost promise, that so, amid the sundry and manifold changes of the world, our hearts may surely there be fixed, where true joys are to be found.
Through our Lord.}


    
    \smalltitle{V Sunday after Easter.}

\texteparacol{
\lettrine{D}{e}us, a quo bona cun\-cta procédunt, largíre supplícibus tuis:~†  ut cogitémus te inspiránte quæ recta sunt;~* et, te gubernánte éadem faciámus. Per Dóminum nostrum.}{english}
{\noindent O God, from whom all good things do come, grant to us thy humble servants that by thy holy inspiration we may think those things that be good, and by thy merciful guiding may perform the same.
Through our Lord.}

\myrule

\smalltitle{Commemoration of the Cross.}

\gscore[Ant.]{6.}{an_crucifixus_solesmes_1961}
\antiphona{english}{He that was crucified is risen from the dead, and hath redee med us. Alleluia, Alleluia}
\smallskip

   \texteparacol{\vv Dícite in natiónibus, allelúia.

\rr Quia Dóminus regnávit a ligno, allelúia.

Orémus.

\lettrine{D}{e}us qui pro nobis Fílium tuum Crucis patíbulum subíre voluísti, ut inimíci a nobis expélleres potestátem: * concéde nobis fámulis tuis; ut resurrectiónis grátiam consequámur. Per eúmdem Dóminum.

\rr Amen.}{english}{
\noindent \vv Say among the heathen, alleluia.

\noindent \rr That the Lord reigneth from the tree, alleluia.

\noindent Let us pray. O God, who didst send thy Son to suffer death for us upon the Cross, that thou mightest deliver us from the power of the enemy grant unto us thy servants to be made partakers of his Resurrection.
Through the same Lord.}

\textes{dv_la_en}

\myrule

\orrubrique

\smallscore{domine_exaudi_solemn}

\myrule

\textes{cantor_rubric}

\gscore[]{7.}{BD_TP}

\noindent \vv Let us bless the Lord.

\noindent \rr Thanks be to God.

\textebilingue{fidelium_animae}{fidelium_animae_en}

\rubrique{\begin{enpars}
Then the Lord's Prayer is recited silently, and the office is concluded with the following verses, the Marian antiphon, etc., unless Compline follows, or, in parish and other churches not bound to the office in choir, Benediction of the Blessed Sacrament or a sermon, or even some other pious exercises.\end{enpars}}

\rubrique{\begin{enpars}The versicle \normaltext{Dóminus} and response \normaltext{Et vitam } are sung on a lower note.\end{enpars}}

\texteparacol{
\noindent \vv Dóminus det nobis suam pacem.

\noindent \rr Et vitam ætérnam. Amen.}
{english}
{\noindent \vv May the Lord grant us his peace.

\noindent \rr And life everlasting. Amen.}

\rubrique{\begin{enpars}In Paschal Time, on Sundays and weekdays alike, all remain standing.\end{enpars}}

\smalltitle{Regina Cæli.}

\rubrique{¶ From Compline of Holy Saturday until None of Saturday after Pentecost inclusive.}

%\gscore[]{6.}{an--regina_caeli--solesmes}
%%\vspace{1ex}
%
%\smalltitle{Simple Tone.}

\gscore[]{6.}{an_regina_caeli_simple_tone_solesmes}

\textes{regina_caeli_la_en}

\rubrique{The versicle and response are sung on a lower note.}
\label{conclusion} \phantomsection

\textes{divinum_auxilium_la_en}

\end{document}